\section{Intended Users of the Product}

The proposed system is intended for academic stakeholders who participate in practical laboratory examinations before, during, and after the scheduled session. Each user group interacts with the platform differently according to its responsibilities, but all of them contribute to a controlled and accountable exam workflow \cite{sommerville}.

\subsection{Administrator}
The administrator uses the system to manage user accounts, assign roles and permissions, register laboratories, map available machines, configure institutional settings, and review overall exam records \cite{rbac1}. This role is responsible for maintaining operational readiness, ensuring that approved users can access the appropriate modules, and generating administrative reports for audit and policy review.

\subsection{Teacher}
Teachers use the system to create exams, upload instructions, schedule lab sessions, review submissions, manage results, and study generated reports after the exam \cite{eassess1}. Teachers are responsible for defining the academic content and operational parameters of each practical examination session.

\subsection{Invigilator}
Invigilators use the live monitoring dashboard during the exam to observe active student sessions, track timing and progress status, respond to alerts, log suspicious incidents, and take corrective action when necessary \cite{mit_cbt}. This role is central to the real-time invigilation support that the system is designed to provide. The invigilator may be the same person as the teacher or may be a separately assigned academic staff member.

\subsection{Student}
Students use the system to log in to their assigned exam session, complete identity verification, access instructions, perform the practical task in a controlled environment, and submit the required files within the scheduled time \cite{biometric1}. Students interact with a focused and time-bound interface that restricts access to only the authorized session.

